\chapter{基数}
\label{chap:cardinality}

正整数与正偶数，哪一个更多？正偶数只是正整数的一部分，似乎更少；
可是给每个正整数$n$配上正偶数$2n$，两边的元素又恰好一一配完。
有理数在数轴上处处稠密，那么实数与有理数能否一一配对？
无限集合无法像有限集合那样逐个数完；不过，有限集合的元素个数也可以用它与$\{1,\ldots,n\}$之间的双射来确定。
本章把“能否一一配对”当作比较大小的标准，推广到任意集合。

本章先用双射定义等势与基数，再用单射比较基数大小，证明两个方向的单射足以推出等势、幂集总严格更大；
最后研究可数集，整理可数性在哪些构造下保持，并证明有理数可数而实数不可数。
例子中用到读者熟悉的数系及其十进制表示，实数的确界与完备性在第~\ref{chap:real-complex-functions}章建立。

\section{等势与基数的定义}
\label{sec:comb-cardinality}
\label{sec:comb-equinumerosity-cardinality}

\paragraph{等势与基数}
对任意两个集合，即使数不出个数，也可以问能否把两边的元素恰好一一配完。若存在双射$f:A\to B$，就称集合$A$与$B$\term{等势}（equinumerous）。

等势具有三条性质：
\begin{enumerate}
\item 恒等映射$\operatorname{id}_A$说明$A$与自身等势（自反）；
\item 双射的逆映射仍是双射，所以$A$与$B$等势时$B$与$A$也等势（对称）；
\item 双射的复合仍是双射，所以$A$与$B$、$B$与$C$等势时$A$与$C$等势（传递）。
\end{enumerate}

等势只看能否配对，不依赖元素的名称：把$\{1,2,3\}$的元素改名为$a,b,c$，改名映射本身就是双射。

等势的集合称为具有相同的\term{基数}（cardinality），记作$|A|=|B|$。
$|A|=|B|$恰好意味着$A$与$B$等势；本章只用这一记号表达比较，不构造基数本身。

\paragraph{有限集合与无限集合的基数}
若$A$与某个$\{1,\ldots,n\}$等势（$n\geq0$），就称$A$为有限集合，它的基数沿用上一章的元素个数，$|A|=n$；
上一章已说明这样的$n$唯一，并约定$|\varnothing|=0$。

与任何$\{1,\ldots,n\}$都不等势的集合称为无限集合。
例如$\mathbb N_+$是无限集合：若有双射$\{1,\ldots,n\}\to\mathbb N_+$，它的$n$个值中最大者加$1$不会被取到，与满射矛盾。

有限集合的真子集元素总是更少，上一章已由加法规则得到$B\subsetneq A$时$|A|>|B|$。无限集合却可能与自己的真子集等势。
映射$n\mapsto2n$从$\mathbb N_+$到正偶数集合是双射：不同的$n$给出不同的$2n$，每个正偶数$2m$都由$m$得到。
所以正偶数集虽是$\mathbb N_+$的真子集，基数却与$\mathbb N_+$相同。
可见“删去一部分元素后必定变少”不能搬到无限集合上，比较无限集合的大小只能依靠映射。

\section{基数的性质}
\label{sec:comb-cardinality-properties}

等势只回答“是否一样多”。本节用单射定义“不超过”，并讨论两个方向的“不超过”能否推出一样多，以及是否存在最大的基数。

\paragraph{基于单射的比较}
若存在单射$f:A\to B$，$A$的每个元素都能占据$B$中一个不同的位置，这时记$|A|\leq|B|$。
若$|A|\leq|B|$但$A$与$B$不等势，记$|A|<|B|$。

双射是单射，所以等势的集合在两个方向上都满足$\leq$。

单射的复合仍是单射，所以$|A|\leq|B|$与$|B|\leq|C|$推出$|A|\leq|C|$，这种比较具有传递性。

若$A\subseteq B$，包含映射$a\mapsto a$是单射，所以子集的基数不超过原集合。

但真包含$A\subsetneq B$并不能推出$|A|<|B|$，正偶数集就是例子：要证明严格不等式，必须排除一切可能的双射。

\paragraph{双向单射推出等势}
若两个方向各有一个单射，两个集合是否一定一样多？对有限集合，比较元素个数即可得到肯定的回答；
对无限集合，两个单射各管一个方向，并不能直接拼成双射，需要另行构造。
\begin{theorem}[Cantor--Schröder--Bernstein定理]
\label{thm:comb-cantor-schroeder-bernstein}
若存在单射$f:A\to B$与$g:B\to A$，则存在双射$h:A\to B$，即$|A|=|B|$。
\end{theorem}
\begin{proof}
我们要为每个$a\in A$选定一个$B$中的配对对象，可选的只有两种：用$f$送到$f(a)$，或者在$a=g(b)$时沿$g$反向找回$b$。
不在$g(B)$中的元素没有反向对象，只能用$f$。设$a$用了$f$，那么$f(a)$已被占用，
元素$g(f(a))$就不能再反向找回$f(a)$，也只能用$f$；如此连锁下去。据此，
令$A_0=A\setminus g(B)$，并递推地令$A_{n+1}=g(f(A_n))$，即对$A_n$先施行$f$再施行$g$；
记$C=\bigcup_{n\geq0}A_n$，它收集了所有“只能用$f$”的元素，其余元素则用$g$的反向对应。定义$h:A\to B$：对$a\in C$令$h(a)=f(a)$；
对$a\in A\setminus C$，令$h(a)$为满足$g(b)=a$的那个$b$。
后一种规定有意义：$A_0\subseteq C$，所以$a\notin C$蕴含$a\notin A_0$，即$a\in g(B)$，这样的$b$存在；
又因为$g$是单射，它唯一。于是在$A\setminus C$上，$h$就是$g$的反向对应，满足$g(h(a))=a$。

下面验证$h$是双射。先看两块的像互不重叠：若某个$a\in C$与$a'\in A\setminus C$有$f(a)=h(a')$，
则$a'=g(h(a'))=g(f(a))$；而$a\in A_n$对某个$n$成立，于是$a'\in g(f(A_n))=A_{n+1}\subseteq C$，与$a'\notin C$矛盾。
再看每一块内部：$f$是单射，所以$h$在$C$上单射；在$A\setminus C$上，$h(a)=h(a')$两边作用$g$得$a=a'$。
两块内部单射且像不重叠，所以$h$是单射。最后看两块的像覆盖$B$：任取$b\in B$，若$g(b)\notin C$，
则$g(b)\in A\setminus C$，$h(g(b))$是满足$g(b')=g(b)$的$b'$，由$g$单射即$b$本身；
若$g(b)\in C$，由于$g(b)\in g(B)$，它不在$A_0$中，于是$g(b)\in A_{n+1}=g(f(A_n))$对某个$n$成立，
即存在$a\in A_n$使$g(b)=g(f(a))$，由$g$单射得$b=f(a)=h(a)$。所以$h$是满射，从而是双射。
\end{proof}
这一定理说明单射比较具有反对称性；连同自反性与传递性，基数的“不超过”在形式上与偏序相同。
证明两个集合等势时，只需在两个方向上各给一个单射，不必直接写出双射。

下面用这一方法证明区间$(0,1)$与实轴$\mathbb R$等势。一个方向，包含映射是$(0,1)$到$\mathbb R$的单射。

另一个方向，取
\[
 \varphi(x)=\frac12+\frac{x}{2(1+|x|)},\qquad x\in\mathbb R,
\]
需要验证两点。
\begin{itemize}
 \item 取值范围：由$|x|<1+|x|$得$\bigl|\frac{x}{2(1+|x|)}\bigr|<\frac12$，所以$\varphi(x)\in(0,1)$。
 \item 单射性：$\varphi(x)-\frac12$与$x$同号，所以$\varphi(x)=\varphi(y)$时$x,y$同号或同为零。
       若$x,y\geq0$，由$\frac{x}{1+x}=\frac{y}{1+y}$交叉相乘得$x=y$；若$x,y<0$，由$\frac{x}{1-x}=\frac{y}{1-y}$同样得$x=y$。
\end{itemize}
于是$\varphi$是$\mathbb R$到$(0,1)$的单射。

由Cantor--Schröder--Bernstein定理，$(0,1)$与$\mathbb R$等势。

长度为$1$的区间与整条实轴等势，可见基数不反映长度。

\paragraph{幂集总有更大的基数}
有限集合满足$|\mathcal P(A)|=2^{|A|}>|A|$。对任意集合，这一严格不等式仍然成立。
\begin{theorem}[Cantor定理]
\label{thm:comb-cantor-power-set}
对任意集合$A$，都有$|A|<|\mathcal P(A)|$。
\end{theorem}
\begin{proof}
映射$a\mapsto\{a\}$把不同的元素送到不同的单元素集，是$A$到$\mathcal P(A)$的单射，所以$|A|\leq|\mathcal P(A)|$。
再设$f:A\to\mathcal P(A)$是任一映射，每个$f(a)$都是$A$的子集，可以问$a$本身是否属于它。构造
\[
 D=\{a\in A:a\notin f(a)\}\in\mathcal P(A).
\]
若$f$是满射，则存在$d\in A$使$f(d)=D$。按$D$的成员条件，$d\in D$当且仅当$d\notin f(d)$，
而$f(d)=D$，这就是说$d\in D$当且仅当$d\notin D$：若$d\in D$，则$d\notin D$；若$d\notin D$，则$d\in D$，两种情形都矛盾。
所以任何$f$都不是满射，更不是双射，结合前面的单射便得$|A|<|\mathcal P(A)|$。
\end{proof}
证明的关键是$D$的构造：对每个$a$，$D$与$f(a)$至少在元素$a$上不同，因为$a$恰好属于二者之一，
所以$D$不可能等于任何$f(a)$。这种逐个对象制造差异的构造称为对角线方法，后面证明实数不可数时还会用到。

上一章用零一记录数出了有限集合的子集个数。对任意集合，同样的对应仍然成立：每个子集$S\subseteq A$可由它的成员记录表示，
即$a\in S$时取$1$、否则取$0$的映射$\chi_S:A\to\{0,1\}$；反过来，每个映射$z:A\to\{0,1\}$都是子集$z^{-1}(\{1\})$的成员记录。
所以$S\mapsto\chi_S$是$\mathcal P(A)$到$\{0,1\}^A$（$A$到$\{0,1\}$的全部映射）的双射，二者等势。
有限情形下$|\{0,1\}^A|=2^{|A|}$；对无限集合，记号$2^{|A|}$就表示$\{0,1\}^A$的基数，也就是幂集的基数。

对任一集合反复取幂集，得到$A,\mathcal P(A),\mathcal P(\mathcal P(A)),\ldots$，基数一个比一个大。
所以不存在最大的基数：任给一个集合，它的幂集都更大。这一结论只讨论任意给定的集合及其幂集，
不把“所有集合”当作一个集合；若把全部集合收集成一个集合$V$，幂集$\mathcal P(V)$的元素都是集合，因而$\mathcal P(V)\subseteq V$，
基数不超过$V$，与Cantor定理矛盾。

\section{可数集}
\label{sec:comb-countable-sets}

能用正整数编号列尽的集合，是后文讨论数列、可数求和与可数覆盖的基础。
本节给出定义与列表刻画，整理可数性的保持规则，再说明有理数可数、实数和无限零一序列不可数。

\subsection{可数集的定义}
\label{sec:countable-set-definition}

\begin{definition}[可数集]
\label{def:analysis-countable-set}
能与正整数集$\mathbb N_+$等势的集合称为可数无限集。
有限集与可数无限集统称\term{可数集}（countable set）；不是可数集的集合称为\term{不可数集}（uncountable set）。
\end{definition}
可数无限集的基数记为$\aleph_0$，读作“阿列夫零”；有限集也算可数集，但其基数是非负整数。
由等势的对称性与传递性，与可数集等势的集合仍可数，与不可数集等势的集合仍不可数。

按定义证明可数，需要写出没有重复的编号；实际构造时，往往只能先写出一个可能有重复的列表。
下面的定理说明，只要列表不遗漏元素，重复并无妨碍。
\begin{theorem}[可数集的列表刻画]
\label{thm:comb-countable-surjection}
非空集合$A$可数，当且仅当存在满射$g:\mathbb N_+\to A$。
\end{theorem}
\begin{proof}
若$A$是可数无限集，取双射$g:\mathbb N_+\to A$即可。若$A=\{a_1,\ldots,a_m\}$有限且$m\geq1$，
令$n\leq m$时$g(n)=a_n$、$n>m$时$g(n)=a_m$，即重复末项把有限列表延长，它是满射。
反之，设$g$是满射。沿$n=1,2,\ldots$依次检查$g(n)$，只保留此前未出现过的值，即按首次出现去重。
每个$a\in A$都等于某个$g(n)$，因而在第$n$步或更早被保留；保留下来的元素互不相同。
若保留的元素只有有限个，$A$是有限集；否则按保留的先后编号$1,2,\ldots$，得到$\mathbb N_+$与$A$的双射。
\end{proof}
于是，要证明一个非空集合可数，只需给出一个可以有重复、但不遗漏任何元素的列表。
空集没有这样的列表，但它是有限集，本身可数，以下各结论中都单独处理。

\subsection{可数集的性质}
\label{sec:countable-set-properties}

\paragraph{子集与函数的像}
由列表刻画，删去一部分元素或用映射把元素送到另一个集合，都不会把可数集变成不可数集。
\begin{theorem}[子集与像的可数性]
\label{thm:comb-countable-subset-image}
设$A$可数。则$A$的每个子集可数；对任意映射$f:A\to Y$，像$f(A)$可数。
特别地，若存在单射$h:B\to\mathbb N_+$，则$B$可数。
\end{theorem}
\begin{proof}
空集以及$A=\varnothing$时的像$f(\varnothing)=\varnothing$都可数，以下设$A$非空，取满射$g:\mathbb N_+\to A$。
对像，逐项取像得到列表$f(g(1)),f(g(2)),\ldots$，即$f\circ g:\mathbb N_+\to f(A)$，它仍是满射，由列表刻画$f(A)$可数。
对非空子集$E\subseteq A$，取定$e_0\in E$，把列表中不属于$E$的项都换成$e_0$，即定义$r(a)=a$（$a\in E$）、$r(a)=e_0$（$a\notin E$）；
$r:A\to E$是满射，所以$E=r(A)$是可数集的像，因而可数。
最后，单射$h$给出$B$与$h(B)$之间的双射，而$h(B)\subseteq\mathbb N_+$可数，故$B$可数。
\end{proof}
例如素数集可数。最后一条是常用判据：给$B$的元素指定互不相同的正整数编码，$B$就可数。

\paragraph{有限笛卡尔积}
怎样列出全部正整数对？若逐行列举$\mathbb N_+^2$，就会在第一行$(1,1),(1,2),\ldots$永远进行下去，到不了第二行。
改按坐标和由小到大分层列举：先$(1,1)$，再$(1,2),(2,1)$，再$(1,3),(2,2),(3,1)$，依此类推。
坐标和为$s$的一层只有$s-1$项，所以每层有限；任意一对$(i,j)$都在坐标和为$i+j$的那一层出现，且只出现一次。
因此这份列表在有限步内到达每一对，给出$\mathbb N_+$与$\mathbb N_+^2$的双射，正整数对的集合可数。
借助这份列表，可以处理任意有限个可数集的乘积。
\begin{theorem}[有限笛卡尔积的可数性]
\label{thm:comb-countable-product}
若$A_1,\ldots,A_k$可数，则$A_1\times\cdots\times A_k$可数。
\end{theorem}
\begin{proof}
若某个$A_i$为空，乘积为空集，结论成立。以下设各$A_i$非空，对$k$归纳，$k=1$时无需证明。
设$P=A_1\times\cdots\times A_k$可数；它非空，可取满射$g:\mathbb N_+\to P$与$g':\mathbb N_+\to A_{k+1}$。
映射$(m,n)\mapsto(g(m),g'(n))$是$\mathbb N_+^2$到$P\times A_{k+1}$的满射，
所以$P\times A_{k+1}$是可数集$\mathbb N_+^2$的像，由定理~\ref{thm:comb-countable-subset-image}可数；
它与$A_1\times\cdots\times A_{k+1}$之间有双射，故后者也可数。
\end{proof}
例如$\mathbb N^d$对每个正整数$d$都可数。

\paragraph{可数个可数集的并}
若有可数多组对象，每组都能列举，能否给合并后的全部对象统一编号？
\begin{theorem}[可数并的可数性]
\label{thm:comb-countable-union}
设$(A_n)_{n\in\mathbb N_+}$中每个$A_n$可数，并且对每个非空的$A_n$已给定一个满射$g_n:\mathbb N_+\to A_n$。
则$\bigcup_{n\in\mathbb N_+}A_n$可数。
\end{theorem}
\begin{proof}
设$I=\{n:A_n\neq\varnothing\}$，若$I$为空则并集为空。否则$I$作为$\mathbb N_+$的子集可数，
$I\times\mathbb N_+$由定理~\ref{thm:comb-countable-product}可数。映射$(n,j)\mapsto g_n(j)$是$I\times\mathbb N_+$到
$\bigcup_nA_n$的满射：并集中的每个元素属于某个非空$A_n$，又是$g_n$的某个值。由定理~\ref{thm:comb-countable-subset-image}，并集可数。
\end{proof}
证明相当于把双索引的全部$g_n(j)$按$(n,j)$的坐标和分层列出，重复出现的元素只保留第一次。
同一元素可以属于多个$A_n$，也可以在同一个列表中多次出现；这些重复只影响编号，不影响可数性。

定理要求事先给定各组的列表，本书遇到的可数并都带有明确的列表。
有限个可数集的并也包括在内，只需令其余的$A_n$为空集。
例如，全部有限长度的自然数序列组成的集合可数：长度为$k\geq1$的序列组成$\mathbb N^k$，由有限积的结论可数，
并且可以按前面的分层方式写出列表；长度为$0$的只有空序列。全部有限长度序列是这些集合对$k=0,1,2,\ldots$的并，所以仍可数。

\paragraph{有理数的可数性}
有理数在数轴上稠密：任意两个不同有理数之间还有无限多个有理数，似乎无法逐个列出。上面的保持规则却说明它们可以列尽。
\begin{theorem}[有理数可数]
\label{thm:card-rationals-countable}
整数集$\mathbb Z$与有理数集$\mathbb Q$都是可数集，$|\mathbb Z|=|\mathbb Q|=\aleph_0$。
\end{theorem}
\begin{proof}
把整数排成$z_1,z_2,z_3,\ldots=0,1,-1,2,-2,\ldots$，即第$2k$项是$k$、第$2k+1$项是$-k$；每个整数恰好出现一次，所以$\mathbb Z$是可数无限集。
每个有理数都能写成$p/q$，其中$p\in\mathbb Z$、$q\in\mathbb N_+$，所以$\mathbb Q$是$\mathbb Z\times\mathbb N_+$在映射$(p,q)\mapsto p/q$下的像。
由有限笛卡尔积与像的可数性，$\mathbb Q$可数；又因为$\mathbb Q\supseteq\mathbb N_+$不是有限集，它是可数无限集。
\end{proof}
证明中的列表可以具体写出：按正整数对$(i,j)$的坐标和分层，依次写出$z_i/j$，得到
\[
 \frac{0}{1},\ \frac{0}{2},\ \frac{1}{1},\ \frac{0}{3},\ \frac{1}{2},\ \frac{-1}{1},\ \ldots.
\]
列表有重复，例如$0/1$与$0/2$表示同一个有理数；按列表刻画，只保留每个有理数首次出现的位置，就得到$\mathbb N_+$与$\mathbb Q$的双射。
同理$\mathbb Q^d$可数，即坐标全为有理数的点在$\mathbb R^d$中只有可数多个。稠密与可数互不妨碍。

\paragraph{实数的不可数性}
\phantomsection
\label{sec:comb-real-uncountable}
有理数能被列尽，实数却不能。证明用到读者已有的十进制表示知识：每个实数都有十进制表示，
同一实数有两种表示只发生在以全$0$与全$9$结尾的情形，例如$0.5000\cdots=0.4999\cdots$；
这些事实依赖实数的完备性（第~\ref{sec:ru-supremum-completeness}节的确界原理与闭区间套定理），本章不展开其证明。
\begin{theorem}[实数不可数]
\label{thm:card-reals-uncountable}
开区间$(0,1)$与实数集$\mathbb R$都是不可数集。
\end{theorem}
\begin{proof}
$(0,1)$含有全部$1/(n+1)$，不是有限集。假设它是可数无限集，其元素可以列为$x_1,x_2,\ldots$。
给每个$x_n$选定一个不以无限多个$9$结尾的十进制表示$x_n=0.d_{n,1}d_{n,2}d_{n,3}\cdots$，
构造$y=0.c_1c_2c_3\cdots$：若$d_{n,n}=1$就取$c_n=2$，否则取$c_n=1$。则$y\in(0,1)$，并且对每个$n$，$y$的第$n$位与$x_n$的第$n$位不同。
$y$只用数字$1,2$，其表示唯一；$x_n$的表示也不以全$9$结尾。所以第$n$位不同意味着$y\neq x_n$，$y$不在列表中，矛盾。
因此$(0,1)$不可数。若$\mathbb R$可数，由子集的可数性，$(0,1)$也可数，矛盾，所以$\mathbb R$不可数。
\end{proof}

实数集的基数记为$\mathfrak c$，称为连续统的基数（cardinality of the continuum）。
对端点$a<b$的任一区间$I$，一方面$I\subseteq\mathbb R$；另一方面线性映射$t\mapsto a+(b-a)t$把$(0,1)$双射到$(a,b)\subseteq I$，
与单射$\varphi:\mathbb R\to(0,1)$复合得到$\mathbb R$到$I$的单射。由Cantor--Schröder--Bernstein定理，$|I|=\mathfrak c$：
长度不同的区间仍一样多。又因为$\mathbb N_+\subseteq\mathbb R$而$\mathbb R$不可数，有$\aleph_0<\mathfrak c$。

\paragraph{幂集与无限序列的不可数性}
\phantomsection
\label{sec:comb-uncountable-power-set}
回看Cantor定理：$\mathcal P(\mathbb N_+)$不是有限集，又由$|\mathbb N_+|<|\mathcal P(\mathbb N_+)|$不与$\mathbb N_+$等势，所以不可数；
经成员记录，全部无限零一序列组成的$\{0,1\}^{\mathbb N_+}$与它等势，也不可数。
这与实数的证明是同一种对角线构造：$D$在元素$a$上与$f(a)$不同，$y$在第$n$位上与$x_n$不同。
$\{0,1\}^{\mathbb N_+}$可以看成可数多个两元集的笛卡尔积，所以有限积、可数并保持可数性，并不意味着无限积也保持可数性。

\section{本章小结}
\label{sec:cardinality-summary}

本章把“能否一一配对”作为比较集合大小的标准：双射定义等势与基数，单射定义“不超过”。
相应地，每一种映射构造给出一种结论。给出双射，就证明了等势；给出单射，就证明了不超过；
两个方向各给出一个单射，由Cantor--Schröder--Bernstein定理也能推出等势。
证明严格更大需要先建立单射，再排除两个集合等势；Cantor定理通过排除反方向的所有满射做到了后一步，由此得出幂集总更大、不存在最大的基数。
无限集合可以与真子集等势，长度不同的区间也可以等势，所以基数不反映次序、位置或长度。

以正整数集为标准，可数集就是能被一个（允许重复的）列表列尽的集合。子集、映射的像、有限笛卡尔积与给定列举的可数并都保持可数性，
因此整数、正整数对、有限长度序列与有理数都可数；无限零一序列、正整数集的幂集与实数集则不可数，$\aleph_0<\mathfrak c$。
有理数稠密而可数，实数却不可数，说明数轴上“处处都有”与“能够编号”是两回事。
实数为什么有十进制表示，确界与极限为什么存在，都要从实数的完备性出发说明。
下一部分从实数的确界与完备性开始，建立数列与函数的极限；后文讨论可数覆盖与可数求和时，还会反复用到本章的可数性规则。
